% BBFlow NSF-style Report --- Times 12pt, single spacing, 1-inch margins
% Project Summary: 1 page | Project Description: 6 pages | References separate
\documentclass[12pt]{article}

\usepackage[T1]{fontenc}
\usepackage{mathptmx}          % Times New Roman
\usepackage[margin=1in]{geometry}
\usepackage{setspace}
\singlespacing
\usepackage[parfill]{parskip}
\setlength{\parskip}{6pt}
\usepackage{titlesec}
\titleformat{\section}{\normalfont\bfseries\large}{}{0em}{}[\titlerule]
\titleformat{\subsection}{\normalfont\bfseries}{}{0em}{}
\titlespacing*{\section}{0pt}{12pt}{4pt}
\titlespacing*{\subsection}{0pt}{8pt}{2pt}

\usepackage{hyperref}
\hypersetup{colorlinks=true, linkcolor=blue, citecolor=blue, urlcolor=blue}
\usepackage[numbers,sort&compress]{natbib}
\usepackage{booktabs}
\usepackage{microtype}
\usepackage{amsmath}
\usepackage{amssymb}
\usepackage{graphicx}

\pagestyle{plain}

\begin{document}

% ============================================================
%  PROJECT SUMMARY  (1 page)
% ============================================================
\begin{center}
  {\large\bfseries Project Summary}\\[4pt]
  \textbf{BBFlow: Learning Conformational Ensembles of Proteins\\
  from Backbone Geometry via SE(3) Flow Matching}
\end{center}

\noindent\textbf{Overview.}
Protein function is governed not by a single static structure but by a
\emph{conformational ensemble} drawn from the Boltzmann distribution. Molecular
dynamics (MD) simulation is the established route to such ensembles, but the
cost of overcoming free-energy barriers and reaching ergodicity makes broad
sampling prohibitive for most proteins. Recent deep generative MD emulators
(AlphaFlow, ESMFlow, ConfDiff, BioEmu) approach this problem by fine-tuning
large pre-trained folding models such as AlphaFold~2 or ESMFold, which
inherit a costly dependence on Multiple Sequence Alignments (MSAs) or protein
language model embeddings -- so-called \emph{evolutionary information}. This
dependence is slow at inference, biases predictions toward natural sequences,
and breaks down on \emph{de novo} designs and intrinsically novel sequences
where evolutionary signal is absent. The proposed work introduces
\textbf{BBFlow}~\cite{wolf2025bbflow}, an SE(3) flow-matching model that
generates protein backbone ensembles conditioned solely on the equilibrium
backbone \emph{geometry} of the target protein, eliminating the dependence on
folding-model weights and evolutionary information altogether.

\noindent\textbf{Intellectual Merit.}
Three specific aims will be pursued. (1) \emph{Reformulate ensemble generation
as structure-conditioned flow matching}: instead of $p(x)$ conditioned on
sequence, learn $p(x \mid x_{\mathrm{eq}})$ on the Riemannian manifold
$\mathrm{SE}(3)^N$, with conditioning provided by the equilibrium backbone via
a binned pairwise-distance map and equivariant local-frame direction vectors.
(2) \emph{Introduce a structure-conditioned prior $p_0(x_0 \mid x_{\mathrm{eq}})$}
through partial geodesic interpolation between unconditional noise and
$x_{\mathrm{eq}}$ -- an idea inspired by stochastic interpolants and partial
denoising, but novel in the flow-matching ensemble setting. (3) \emph{Validate
on the ATLAS MD benchmark, on a curated set of \emph{de novo} proteins, and on
multi-chain complexes}, comparing against AlphaFlow, ESMFlow-T, ConfDiff,
BioEmu, and distilled variants of AlphaFlow on RMSF, pairwise RMSD, PCA
$\mathcal{W}_2$, DCCM, transient-contact Jaccard, and per-conformation inference
time. Preliminary published data already establish feasibility: BBFlow trained
from scratch in three GPU-days reaches accuracy competitive with AlphaFlow-T
while being roughly $40\times$ faster, generalizes to \emph{de novo} proteins
where MSA-dependent baselines collapse, and -- although trained only on
monomers -- transfers to multi-chain assemblies. The conditional prior, the
distance encoding, and the index-free architecture are each shown to be
necessary in ablation.

\noindent\textbf{Broader Impacts.}
Decoupling MD emulation from large folding models lowers the compute cost of
ensemble generation by more than an order of magnitude, opening MD-quality
sampling to laboratories and pipelines without TPU-scale resources. Public
release of model weights, training and inference code, and benchmark scripts
(\url{https://github.com/graeter-group/bbflow}) supports reproducibility and
reuse. By extending ensemble prediction to \emph{de novo} and multi-chain
proteins -- regimes uncovered by current state-of-the-art models -- BBFlow
enables \emph{flexibility-aware protein design}, allowing dynamics screening to
be incorporated directly into design loops for binders, biocatalysts, and
allosteric switches. Trainees gain experience in geometric deep learning,
Riemannian generative modelling, and quantitative biophysics, and the work
provides instructional examples in NeurIPS-level ML and structural biology
courses.

\textbf{Keywords:} flow matching; SE(3) equivariance; conformational ensemble;
molecular dynamics emulation; protein backbone; de novo proteins.

\newpage

% ============================================================
%  PROJECT DESCRIPTION  (6 pages)
% ============================================================
\begin{center}
  {\large\bfseries Project Description}
\end{center}

\section{Overview, Hypothesis, and Specific Aims}

\noindent\textbf{Objective.} The objective of the proposed work is to develop and
benchmark \textbf{BBFlow}, an SE(3) flow-matching generative model that produces
short-timescale, MD-consistent backbone conformational ensembles of proteins
conditioned solely on the protein's equilibrium backbone geometry, without any
reliance on multiple-sequence-alignment (MSA) features, protein language model
embeddings, or pre-trained folding-model weights.

\noindent\textbf{Central Hypothesis.} The pairwise-distance map and local-frame
directional encoding of an equilibrium backbone $x_{\mathrm{eq}}$ contain
sufficient geometric information to determine, up to MD precision, the
short-timescale Boltzmann distribution $p(x \mid x_{\mathrm{eq}})$ of backbone
conformations -- and conditioning a flow-matching prior on $x_{\mathrm{eq}}$
yields a model that is faster, more transferable, and less biased than
sequence-conditioned ensemble generators built atop large folding models.

\noindent\textbf{Specific Aims.}
\begin{itemize}
  \item \textbf{Aim~1 -- Architecture and conditioning.} Adapt the GAFL/FrameDiff
  SE(3)-equivariant graph network with Clifford Frame Attention (CFA), and
  condition it on $x_{\mathrm{eq}}$ via (i) a binned $C_\alpha$ pairwise distance
  map and (ii) equivariant pairwise direction vectors expressed in residue-local
  frames. Remove the absolute residue-index feature to discourage memorization
  and enable transfer to multi-chain inputs.
  \item \textbf{Aim~2 -- Conditional prior on $\mathrm{SE}(3)^N$.} Replace the
  unconditional Gaussian/uniform prior of standard flow matching with a
  \emph{conditional} prior $p_0(x_0 \mid x_{\mathrm{eq}})$ constructed by partial
  geodesic interpolation, $x_0 = \gamma(x_{\mathrm{uncond}}, x_{\mathrm{eq}}, \xi)$,
  with $\xi \in (0,1)$. Train end-to-end with the SE(3) flow-matching loss on the
  ATLAS dataset.
  \item \textbf{Aim~3 -- Benchmarking and transferability.} Evaluate on the
  ATLAS test set, on a held-out set of \emph{de novo} proteins generated by
  RFdiffusion and FrameFlow, and on five multi-chain protein complexes. Compare
  against AlphaFlow, AlphaFlow-T, ESMFlow-T, ConfDiff, BioEmu, and distilled
  AlphaFlow variants on six standard ensemble metrics plus per-conformation
  inference time, with ablations isolating the contribution of each design
  choice.
\end{itemize}

\section{Significance and Background}

\subsection{Conformational ensembles are essential to protein function}
Enzymatic catalysis~\cite{benkovic2008}, allosteric regulation~\cite{bernetti2024},
and ligand binding~\cite{karplus2005} are all governed by transitions across the
free-energy landscape rather than by a single equilibrium structure. Sampling
this landscape with all-atom MD~\cite{abraham2015gromacs,adcock2006md}, while
physically rigorous, is computationally limited: meaningful sampling for a
single mid-sized protein typically demands hundreds of nanoseconds to
microseconds of trajectory, a wall-clock cost that scales poorly to proteome-
or design-library-scale screens.

\subsection{Folding models predict structures, not ensembles}
The accuracy revolution of AlphaFold~2~\cite{jumper2021af2}, RoseTTAFold-3T~\cite{baek2021rosettafold},
and the protein language model ESMFold~\cite{lin2023esm2} predicts a single
\emph{equilibrium} structure per sequence and gives no direct access to
conformational distributions. Their predictive power derives from
\emph{evolutionary information} encoded either in the MSA processed by the
Evoformer or implicitly in pLM weights -- features that are expensive to compute
and unavailable for genuinely novel (e.g.\ \emph{de novo} designed) sequences.

\subsection{Generative MD emulators and their limitations}
A first generation of generative ensemble models adapts these folding models:
AlphaFlow~\cite{jing2024alphaflow} fine-tunes AlphaFold~2 to emulate the ATLAS
MD dataset~\cite{vandermeersche2024atlas} via flow matching; ESMFlow-T
substitutes the ESMFold backbone; ConfDiff~\cite{wang2024confdiff} replaces the
structure module with a diffusion model; BioEmu~\cite{lewis2025bioemu} targets
broader equilibrium ensembles via a separately trained model. All four inherit
two costs from their parents: (i) per-sample inference of an Evoformer-class
network, yielding $\mathcal{O}(10\text{--}30\,\mathrm{s})$ per conformation, and
(ii) sensitivity to MSA quality, which is structurally absent for \emph{de novo}
proteins. Empirically, AlphaFlow-T over-stabilizes ensembles (median RMSF
$1.17~$\AA{} vs.\ MD's $1.48~$\AA), AlphaFlow without templates fails on
\emph{de novo} proteins, and BioEmu does not target ATLAS-style MD emulation.

\subsection{Significance of the proposed work}
The proposed BBFlow model addresses each limitation in a single architectural
move: by conditioning on \emph{geometry} ($x_{\mathrm{eq}}$) instead of
\emph{sequence}, it removes the Evoformer/pLM bottleneck, runs in
$\sim 0.8~\mathrm{s}$ per conformation, and is independent of evolutionary
information -- the property that makes AlphaFlow brittle on \emph{de novo}
sequences. As MD itself begins from an initial structure, conditioning on
$x_{\mathrm{eq}}$ does not impose any cost above the existing MD workflow, and a
folding-model$+$BBFlow pipeline remains $\sim 30\times$ faster than AlphaFlow even
when only a sequence is available. The work thereby reframes MD emulation as a
\emph{structure-to-ensemble} problem, with direct relevance to dynamics-aware
\emph{de novo} design~\cite{viliuga2025flexibility, komorowska2025nmatune,
guo2025dynamic}.

\section{Innovation}

\subsection{Geometric, evolution-free conditioning}
Prior MD emulators inject conditioning via the Evoformer's processing of an MSA
or via pLM embeddings. The proposed work replaces this with two purely
geometric features computed from $x_{\mathrm{eq}}$:
\begin{align}
  s_{ij} &= \mathrm{bin}\!\left(\|z_i - z_j\|_2\right) \quad\text{(22 bins, $0$--$20\,\text{\AA}$)},\\
  e_{ij} &= r_i^{-1}\!\left(\frac{z_j-z_i}{\|z_j-z_i\|_2}\right), \quad \|z_i-z_j\| < 5\,\text{\AA},
\end{align}
where $(r_i, z_i) \in \mathrm{SE}(3)$ is the local frame of residue~$i$. The
first is a soft contact map -- shown in MSA-based methods to be the implicit
target representation of evolutionary information~\cite{lin2023esm2} -- and the
second is a strictly equivariant directional cue absent from prior ensemble
generators.

\subsection{A conditional prior for flow matching}
Standard flow matching~\cite{lipman2023flowmatching} uses an unconditional prior
$p_0$ (Gaussian translations, uniform $\mathrm{SO}(3)$ rotations). The proposed
work introduces a \emph{conditional} prior obtained by partial geodesic
interpolation between unconditional noise and $x_{\mathrm{eq}}$,
\begin{align}
  x_{\mathrm{uncond}} &\sim p_{\mathrm{uncond}}, \qquad
  x_0 = \gamma(x_{\mathrm{uncond}}, x_{\mathrm{eq}}, \xi), \quad \xi=0.2,
\end{align}
where $\gamma(\cdot, \cdot, 0) = x_{\mathrm{uncond}}$ and $\gamma(\cdot, \cdot, 1)
= x_{\mathrm{eq}}$. Conceptually this is a flow-matching analogue of partial
denoising in diffusion~\cite{lu2024str2str} and a generalization of stochastic
interpolants~\cite{albergo2022interpolants} to manifold-valued conditioning. To
the best of our knowledge, BBFlow is the first ensemble model on $\mathrm{SE}(3)^N$
to adopt a structure-conditioned prior of this form.

\subsection{Index-free, transferable architecture}
Folding-style architectures~\cite{jumper2021af2,jing2024alphaflow,yim2023fast}
embed an absolute residue index, which encodes a fixed chain ordering and ties
the model to single-chain inputs. The proposed architecture removes the index
encoding altogether, relying only on geometric features. Ablation
(Sec.~\ref{sec:approach}) demonstrates that this choice both improves accuracy
and -- more importantly -- enables zero-shot transfer to multi-chain complexes,
which existing baselines cannot ingest without re-engineering.

\subsection{Distinction from contemporaneous work}
ESMFlow-T, AlphaFlow-T, AlphaFlow-T$_{12\mathrm{L,dist}}$ and Li et al.~\cite{li2024improving}
attack the inference-time problem by \emph{compressing} pre-trained folding
models. The proposed work attacks it by \emph{redefining the input
representation}, eliminating the folding-model dependency rather than
distilling it. MDGen~\cite{jing2024mdgen} models full all-atom trajectories
with consistent time evolution; the proposed work trades temporal consistency
for a stationary distribution and gains substantial speed and robustness.

\section{Approach}\label{sec:approach}

\subsection{Aim 1: Geometric conditioning of an SE(3) flow}

\paragraph{Manifold and vector field.}
Each backbone residue is represented by a frame $x_i = (r_i, z_i) \in
\mathrm{SE}(3)$ with $r_i \in \mathrm{SO}(3)$ and $z_i \in \mathbb{R}^3$, and a
protein of length~$N$ is an element of $\mathcal{M} \equiv \mathrm{SE}(3)^N$. A
time-dependent conditional vector field
\begin{align}
  v(x, t, x_{\mathrm{eq}}) : \mathcal{M} \times [0,1] \times \mathcal{M}_{\mathrm{eq}} \to T_x\mathcal{M}
\end{align}
defines the flow $\phi_t$ via $\frac{d}{dt}\phi_t = v(\phi_t, t, x_{\mathrm{eq}})$,
$\phi_0 = x$. Following~\cite{yim2023se3,bose2024se3}, the SE(3) target vector
field decomposes into rotational and translational components,
\begin{align}
  v_{\mathrm{SO}(3)}(r_t, t \mid r_1) = \frac{\log_{r_t}(r_1)}{1 - t}, \qquad
  v_{\mathbb{R}^3}(z_t, t \mid z_1) = \frac{z_1 - z_t}{1 - t}.
\end{align}

\paragraph{Network.}
The vector field is parameterized by an SE(3)-equivariant graph neural network
adapted from GAFL~\cite{wagner2024gafl}, comprising 6 message-passing blocks
of Clifford Frame Attention. Inputs are (i) the noised frames $x_t$ at time~$t$,
(ii) the binned distance map $s_{ij}$, (iii) the local-frame direction vectors
$e_{ij}$ for residue pairs within $5~$\AA, and (iv) a 128-dimensional embedding
of the amino-acid type. Crucially, no absolute residue index is supplied -- the
ordering is recovered from $x_{\mathrm{eq}}$.

\subsection{Aim 2: Conditional prior and training}
The conditional prior $p_0(\cdot \mid x_{\mathrm{eq}})$ is constructed by
geodesic interpolation, with $\xi = 0.2$ chosen by validation. Training
minimizes the SE(3) flow-matching loss
\begin{align}
  \mathcal{L}_{\mathrm{FM}}
  = \mathbb{E}\!\left[\left\| v - \hat{v}(x_t, t, x_{\mathrm{eq}})\right\|^2_{\mathrm{SE}(3)}\right],
  \quad \|v\|^2_{\mathrm{SE}(3)} \equiv \tfrac{1}{2}\,\mathrm{Tr}(v_r v_r^T) + \|v_z\|_2^2,
\end{align}
with $t \sim \mathcal{U}(0,1)$, $(x_1, x_{\mathrm{eq}}) \sim p_{\mathrm{data}}$,
and $x_0 \sim p_0(\cdot \mid x_{\mathrm{eq}})$. An auxiliary loss following
\cite{yim2023se3} stabilizes early training. We train on the
ATLAS~\cite{vandermeersche2024atlas} split of \cite{jing2024alphaflow}
(1265/39/82 train/val/test) for 3 days on two NVIDIA A100-40GB GPUs, using 20
flow timesteps and 250 sampled conformations per protein at inference.

\subsection{Aim 3: Benchmarking}

\paragraph{Metrics.}
We report (i) per-residue Root Mean Square Fluctuation (RMSF) Pearson~$r$ and
mean absolute error (MAE), (ii) median RMSF (probing systematic
over/under-stabilization), (iii) pairwise RMSD MAE, (iv) PCA
$\mathcal{W}_2$-distance against MD on the first two principal components,
(v) Pearson~$r$ of the Dynamic Cross-Correlation Matrix (DCCM), (vi) transient-
contact Jaccard $J_{\mathrm{tr}}$, and (vii) inference time per generated
conformation on a 302-residue test protein.

\paragraph{Preliminary results -- ATLAS benchmark.}
Table~\ref{tab:atlas} summarises performance on the ATLAS test set. BBFlow is
\textbf{$\sim 40\times$ faster} than AlphaFlow-T at comparable accuracy, achieves
the best pairwise RMSD MAE (0.77~\AA), and matches the MD median RMSF (1.49~\AA{}
vs.\ MD 1.48~\AA) -- avoiding the over-stabilization of AlphaFlow-T (1.17~\AA).

\begin{table}[h]
\centering
\small
\begin{tabular}{lcccccc}
\toprule
Model & RMSF $r\uparrow$ & RMSF MAE$\downarrow$ & Pw-RMSD MAE$\downarrow$ & DCCM $r\uparrow$ & PCA $\mathcal{W}_2\downarrow$ & Time [s]$\downarrow$\\
\midrule
AlphaFlow         & 0.86 & 0.59 & 1.35 & 0.86 & 1.47 & 32.0 \\
ConfDiff          & 0.88 & 0.62 & 1.45 & 0.86 & 1.41 & 20.2 \\
AlphaFlow-T       & 0.92 & 0.41 & \textbf{0.91} & 0.89 & 1.28 & 32.6 \\
ESMFlow-T         & 0.92 & 0.52 & 0.94 & 0.89 & 1.48 & 11.2 \\
AlphaFlow-T$_{\mathrm{dist}}$ & 0.92 & 0.68 & 1.41 & 0.88 & 1.43 & 3.3 \\
\textbf{BBFlow (proposed)} & 0.90 & \underline{0.42} & \underline{0.77} & 0.87 & \underline{1.33} & \textbf{0.8} \\
\bottomrule
\end{tabular}
\caption{Performance on the ATLAS test set. BBFlow matches or improves on the
distilled AlphaFlow variants while being more than an order of magnitude
faster than AlphaFlow-T at comparable accuracy. Adapted from~\cite{wolf2025bbflow}.}
\label{tab:atlas}
\end{table}

\paragraph{Preliminary results -- de novo proteins.}
On 50 \emph{de novo} proteins generated by RFdiffusion~\cite{watson2023rfdiffusion}
and FrameFlow~\cite{yim2024framefw} with three 100-ns MD trajectories each,
AlphaFlow without templates collapses to RMSF $r=0.47$ and BioEmu to $r=0.60$,
whereas BBFlow reaches $r=0.84$ with a pairwise RMSD MAE of 0.32~\AA{} -- the
strongest result among methods that do not consume the equilibrium structure
as a template.

\paragraph{Preliminary results -- multi-chain proteins.}
On five protein dimers (1BGS, 8HXR, 5EY7, and two further complexes), DCCM
correlations of 0.85, 0.72, and 0.89 are obtained between BBFlow- and
MD-derived ensembles, capturing both intra- and inter-chain motions despite
training only on monomers. AlphaFlow fails on the same inputs without parser
modifications.

\paragraph{Ablation.}
A targeted ablation (Tab.~\ref{tab:abl}) confirms that (i) the conditional
prior, (ii) the direction encoding, and (iii) removal of the residue-index
feature each contribute; the binned distance encoding is essential -- removing
it raises RMSF MAE from $0.42$ to $5.88$~\AA. Even without amino-acid identity,
BBFlow remains competitive with non-template AlphaFlow.

\begin{table}[h]
\centering
\small
\begin{tabular}{lcccccccc}
\toprule
Variant & Cond.\ prior & Dist.\ enc. & Dir.\ enc. & AA enc. & Index enc. & RMSF MAE & Pw-RMSD MAE & DCCM $r$\\
\midrule
\textbf{BBFlow (full)} & \checkmark & \checkmark & \checkmark & \checkmark &  & \textbf{0.42} & \textbf{0.77} & \underline{0.87}\\
$-$cond.\ prior   &              & \checkmark & \checkmark & \checkmark & \checkmark & 0.52 & 1.15 & 0.85\\
$-$direction enc. &              & \checkmark &              & \checkmark & \checkmark & 0.48 & 0.90 & 0.86\\
$+$index enc.     & \checkmark   & \checkmark & \checkmark & \checkmark & \checkmark & 0.42 & 0.82 & 0.88\\
$-$AA enc.        & \checkmark   & \checkmark &              &              & \checkmark & 0.54 & 0.93 & 0.85\\
$-$distance enc.  & \checkmark   &              &              & \checkmark & \checkmark & 5.88 & 7.08 & 0.55\\
\bottomrule
\end{tabular}
\caption{Ablation on the ATLAS test set. The binned distance encoding is
essential; the conditional prior, direction encoding, and absence of index
encoding each yield measurable improvements. Adapted from~\cite{wolf2025bbflow}.}
\label{tab:abl}
\end{table}

\subsection{Outlook and alternative approaches}

\paragraph{Next steps.}
Three follow-on directions extend the proposed work. (a) \emph{Side-chain and
all-atom ensembles}: BBFlow currently generates backbone frames only; coupling
to a fixed-backbone packer or to MDGen-style atomistic generators is a natural
extension for ligand-binding studies. (b) \emph{Longer-timescale ensembles}:
training on $\mu$s-scale trajectories from D.~E.~Shaw or AlphaFlow-MD would
expand BBFlow beyond the 300-ns regime captured by ATLAS. (c) \emph{Integration
with de novo design}: coupling BBFlow with RFdiffusion in a closed loop would
allow direct optimization of designed flexibility, building on
NMA-tune~\cite{komorowska2025nmatune} and FlexFlow~\cite{viliuga2025flexibility}.

\paragraph{Risks and alternatives.}
\emph{Risk 1: rare events and transient contacts.} BBFlow's transient-contact
Jaccard ($29\%$) trails AlphaFlow-T ($47\%$), suggesting that rare specific
contacts may require evolutionary priors. \emph{Mitigation}: introduce an
optional pLM channel, fine-tune on contact-rich trajectories, or augment the
loss with a contact-prediction term. \emph{Risk 2: dependence on the input
structure.} If $x_{\mathrm{eq}}$ is unavailable, a single AlphaFold~2 prediction
suffices and the AlphaFold~$+$~BBFlow pipeline remains $\sim 30\times$ faster than
AlphaFlow at comparable accuracy. Robustness to noisy inputs is verified by
adding $0.1$--$0.2~$\AA{} Gaussian noise to backbone coordinates and to using random
MD frames as $x_{\mathrm{eq}}$, with only minor degradation. \emph{Risk 3:
distributional scope.} As an MD emulator, BBFlow cannot predict alternative
folded states far from $x_{\mathrm{eq}}$ without retraining on transition-rich
trajectories; specialized models such as Str2Str~\cite{lu2024str2str} can be
used in that regime, accepting reduced MD-emulation fidelity.

\paragraph{Deliverables.}
(D1) Trained BBFlow weights, training code, and evaluation scripts released
under a permissive licence at \url{https://github.com/graeter-group/bbflow};
(D2) public benchmark on ATLAS, \emph{de novo}, and multi-chain test sets;
(D3) a peer-reviewed publication (NeurIPS~2025); (D4) tutorial materials for
incorporating BBFlow into protein design pipelines.

% ============================================================
%  REFERENCES
% ============================================================
\renewcommand{\refname}{\large References}
\begin{thebibliography}{99}\setlength{\itemsep}{2pt}\small

\bibitem{wolf2025bbflow} N.~Wolf, L.~Seute, V.~Viliuga, S.~Wagner, J.~St\"uhmer, F.~Gr\"ater. Learning conformational ensembles of proteins based on backbone geometry. \emph{NeurIPS}, 2025. arXiv:2503.05738.

\bibitem{benkovic2008} S.~J.~Benkovic, G.~G.~Hammes, S.~Hammes-Schiffer. Free-energy landscape of enzyme catalysis. \emph{Biochemistry}, 47(11):3317--3321, 2008.

\bibitem{bernetti2024} M.~Bernetti, S.~Bosio, V.~Bresciani, F.~Falchi, M.~Masetti et al. Probing allosteric communication with combined molecular dynamics simulations and network analysis. \emph{Curr.\ Opin.\ Struct.\ Biol.}, 86:1--10, 2024.

\bibitem{karplus2005} M.~Karplus, J.~Kuriyan. Molecular dynamics and protein function. \emph{PNAS}, 102(19):6679--6685, 2005.

\bibitem{abraham2015gromacs} M.~J.~Abraham et al. GROMACS: high performance molecular simulations through multi-level parallelism from laptops to supercomputers. \emph{SoftwareX}, 1--2:19--25, 2015.

\bibitem{adcock2006md} S.~A.~Adcock, J.~A.~McCammon. Molecular dynamics: survey of methods for simulating the activity of proteins. \emph{Chem.\ Rev.}, 106(5):1589--1615, 2006.

\bibitem{jumper2021af2} J.~Jumper et al. Highly accurate protein structure prediction with AlphaFold. \emph{Nature}, 596:583--589, 2021.

\bibitem{baek2021rosettafold} M.~Baek et al. Accurate prediction of protein structures and interactions using a three-track neural network. \emph{Science}, 373(6557):871--876, 2021.

\bibitem{lin2023esm2} Z.~Lin et al. Evolutionary-scale prediction of atomic-level protein structure with a language model. \emph{Science}, 379(6637):1123--1130, 2023.

\bibitem{jing2024alphaflow} B.~Jing, B.~Berger, T.~Jaakkola. AlphaFold meets flow matching for generating protein ensembles. \emph{ICML}, 2024.

\bibitem{vandermeersche2024atlas} Y.~Vander Meersche, G.~Cretin, A.~Gheeraert, J.-C.~Gelly, T.~Galochkina. ATLAS: protein flexibility description from atomistic molecular dynamics simulations. \emph{Nucleic Acids Res.}, 52(D1):D384--D392, 2024.

\bibitem{wang2024confdiff} Y.~Wang et al. Protein conformation generation via force-guided SE(3) diffusion models. \emph{ICML}, 2024.

\bibitem{lewis2025bioemu} S.~Lewis et al. Scalable emulation of protein equilibrium ensembles with generative deep learning. \emph{Science}, 389(6761):eadv9817, 2025.

\bibitem{lipman2023flowmatching} Y.~Lipman, R.~T.~Q.~Chen, H.~Ben-Hamu, M.~Nickel, M.~Le. Flow matching for generative modeling. \emph{ICLR}, 2023.

\bibitem{albergo2022interpolants} M.~S.~Albergo, E.~Vanden-Eijnden. Building normalizing flows with stochastic interpolants. \emph{ICLR}, 2022.

\bibitem{lu2024str2str} J.~Lu, B.~Zhong, Z.~Zhang, J.~Tang. Str2Str: a score-based framework for zero-shot protein conformation sampling. \emph{ICLR}, 2024.

\bibitem{yim2023se3} J.~Yim et al. Fast protein backbone generation with SE(3) flow matching. arXiv:2310.05297, 2023.

\bibitem{yim2023fast} J.~Yim et al. SE(3) diffusion model with application to protein backbone generation. \emph{ICML}, 2023.

\bibitem{yim2024framefw} J.~Yim et al. Improved motif-scaffolding with SE(3) flow matching. \emph{TMLR}, 2024.

\bibitem{bose2024se3} J.~Bose et al. SE(3)-stochastic flow matching for protein backbone generation. \emph{ICLR}, 2024.

\bibitem{wagner2024gafl} S.~Wagner, L.~Seute, V.~Viliuga, N.~Wolf, F.~Gr\"ater, J.~Stuehmer. Generating highly designable proteins with geometric algebra flow matching. \emph{NeurIPS}, 2024.

\bibitem{watson2023rfdiffusion} J.~L.~Watson et al. De novo design of protein structure and function with RFdiffusion. \emph{Nature}, 2023.

\bibitem{li2024improving} S.~Li et al. Improving AlphaFlow for efficient protein ensembles generation. arXiv:2407.12053, 2024.

\bibitem{jing2024mdgen} B.~Jing, H.~Stark, T.~Jaakkola, B.~Berger. Generative modeling of molecular dynamics trajectories. \emph{NeurIPS}, 2024.

\bibitem{viliuga2025flexibility} V.~Viliuga, L.~Seute, N.~Wolf, S.~Wagner, A.~Elofsson, J.~St\"uhmer, F.~Gr\"ater. Flexibility-conditioned protein structure design with flow matching. \emph{ICML}, 2025.

\bibitem{komorowska2025nmatune} U.~J.~Komorowska, F.~Vargas, A.~Rondina, P.~Lio, M.~Jamnik. NMA-tune: generating highly designable and dynamics aware protein backbones. \emph{ICML}, 2025.

\bibitem{guo2025dynamic} A.~B.~Guo et al. Deep learning-guided design of dynamic proteins. \emph{Science}, 388(6749):eadr7094, 2025.

\end{thebibliography}

\end{document}
